%% ma: L12-B15-0021
%% lop: 12
%% chuong: 5
%% bai: 15
%% dang: 12.15.17
%% loai: TN4
%% muc: VD
%% dap_an: "A"
%% nguon: "(NBV)-10 ĐỀ KIỂM TRA HK2-TOÁN 12-DUNG CHUNG 3 SÁCH.pdf (H03, Đề số 3), Phần 1 Câu 3"
%% kien_thuc: "Bài 14 (vectơ pháp tuyến, vị trí tương đối đường thẳng và mặt phẳng), tích có hướng (Bài 8), Bài 15 (phương trình tham số)"
%% trang_thai: cho-duyet
%% ly_do: "Dạng toán mới đề xuất 12.15.17, chờ giáo viên duyệt"
%% ghi_chu: "Đáp án gốc A đúng."
\begin{ex}
Trong không gian $Oxyz$, cho điểm $A(1;-2;3)$ và hai mặt phẳng $(P)\colon x+y+z+1=0$, $(Q)\colon x-y+z-2=0$. Phương trình nào dưới đây là phương trình đường thẳng đi qua $A$, song song với cả hai mặt phẳng $(P)$ và $(Q)$?
\choice{\True $\begin{cases}x=1+t\\y=-2\\z=3-t\end{cases}$}{$\begin{cases}x=-1+t\\y=2\\z=-3-t\end{cases}$}{$\begin{cases}x=1+2t\\y=-2\\z=3+2t\end{cases}$}{$\begin{cases}x=1\\y=-2\\z=3-2t\end{cases}$}
\loigiai{Hai mặt phẳng có vectơ pháp tuyến $\vec n_P=(1;1;1)$, $\vec n_Q=(1;-1;1)$. Đường thẳng cần tìm có vectơ chỉ phương $\vec u=\left[\vec n_P,\vec n_Q\right]=(2;0;-2)$, chọn $(1;0;-1)$. Đường thẳng qua $A(1;-2;3)$ nên có phương trình $x=1+t$, $y=-2$, $z=3-t$. Chọn A.

Phương án B có cùng phương nhưng không đi qua $A$ (với $y=2\neq-2$).}
\end{ex}
